\scanpage{87}
% archive sheet number: 89

\uline{NB} On sait que la donnée d'un octaèdre \inserted{comb.} (pas orienté) revient à celle d'un système $(I,\omega)\ \deleted{\unsure{\text{dans}}}\ \mathbb{E}\mathrm{ns}_4\times\mathbb{E}\mathrm{ns}_2$, ens des codiagonales et des orientations. Ici $\omega=\{i,-i\}$, \deleted{\unsure{d'autre}} \deleted{part, la donnée de $I$ équi}\inserted{Mais ici ce n'est pas $I$} qui est donné directement, mais $S$, qui est un ens de \deleted{son} cardinal 6 sur lequel $V$ opère.

La condition à mettre \deleted{et} pour que ceci définisse une \deleted{str} structure octaèdrale (dont $S$ sera l'ens des sommets) est que $\forall\sigma\in V^*$, $\deleted{\unsure{V^\sigma|S^\sigma}}\ \operatorname{card}S^\sigma=2$, et $S=\bigcup_{\sigma\in V^*}S^\sigma$ (réunion forcément disjointe) (il s'ensuit que $\sigma$ opère sur $(S\setminus S^\sigma)$ par transposition sur \deleted{des}\inserted{chacune} deux composantes $S^{\sigma'},S^{\sigma''}$ (où $V^*=\{\sigma,\sigma',\sigma''\}$). Donc la donnée de $\deleted{\unsure{S,V\subset|S,V\hookrightarrow}}$ $(V,S,u:V\hookrightarrow\unsure{\mathfrak S|\mathscr S}_S)$, ou $(S,V\subset\mathfrak S_S)$, équivaut à celle de $S$, avec une involution sans pts fixes $a$ (ayant les $S^\sigma$ comme orbites).

Revenant à la donnée d'un $V$ qui opère \inserted{fid\textsuperscript{t}} sur $X$, considérons une orbite \inserted{$I$} de $V$ sur $X^*=X\setminus S(V)$. Donc $I$ est un torseur sous $V$. La donnée de $(X,I,I\hookrightarrow X)$ équivaut à celle de $(X,I,V(I)\hookrightarrow G_X,\text{ orbite de }V(I)\text{ dans }X)$, et on se gardera de confondre le $(V=V(I))$-torseur $I$, avec le $V$-torseur $I(V)\inserted{=I(V,X)}$ défini via l'octaèdre $\mathrm{Oc}(V)=\mathrm{Oc}(V,X)$.

En termes de $\inserted{I\deleted{\unsure{V}},\{\pm i\}}$, on peut définir un octaèdre \deleted{\uline{orienté}} comme un \deleted{\unsure{ens bine univoquement défini par la}} condition que le \deleted{groupe de ses aut.} \inserted{``l'ens de ses \unsure{4 faces|faces}'' soit ``$I$'', l'ens de ses \deleted{``pôles'' aux} orientations $\{\pm i\}$} \deleted{bi-transp} \deleted{soit $V$.} Cet \deleted{octa} octaèdre, \deleted{une fois choisie une $i=\sqrt{-1}$,} peut se réaliser canoniquement \deleted{\unsure{par des som}} comme un $(X,S)$ où $(X(I),S(I))$. On a un isom can
